\documentclass[../../../../main.tex]{subfiles}
\begin{document}

\begin{flushright}
\footnotesize\textit{【自動文字起こし・要確認】}
\end{flushright}

{\bf [解]}
$f(x)=0$ の2解 $\alpha,\beta$ として，$f(x)=(x-\alpha)(x-\beta)$ とおく．
\begin{align*}
f(f(x))=(f(x)-\alpha)(f(x)-\beta)=0 \quad \cdots \text{①}
\end{align*}
が重解 $r$ を持つ．$\alpha\ne\beta$ だから対称性から，$f(x)-\alpha=0$ が重解を持つ時のみ考えれば良い．$f(x)-\alpha=x^2+2x+\alpha\beta-\alpha=0$ の判別式を $D$ として $D=0$ である．
\begin{align*}
\frac{D}{4}=1-(\alpha\beta-\alpha)=0 \quad\therefore\ 1+\alpha-\alpha\beta=0 \quad \cdots \text{②}
\end{align*}
$\alpha+\beta=-2$ $\therefore \beta=-(\alpha+2)$ だから，②に代入して
\begin{align*}
1+\alpha+\alpha(\alpha+2)=0 \quad\Longrightarrow\quad \alpha^2+3\alpha+1=0 \quad\Longrightarrow\quad \alpha=\frac{-3\pm\sqrt5}{2}
\end{align*}
したがって，$(\alpha,\beta)=\left(\dfrac{-3\pm\sqrt5}{2},-\dfrac{1\pm\sqrt5}{2}\right)$（複号同順）となり，
\begin{align*}
a=\alpha\beta=\frac{-1+\sqrt5}{2},\ \frac{-1-\sqrt5}{2}
\end{align*}
である．この時 $r=-1$ となる．


\newpage

\end{document}
